\chapter{Bone Scan}

\section*{What Is This Test?}
A bone scan uses a radioactive tracer that accumulates in areas of high bone turnover. It detects fractures, infections, arthritis, and cancer spread to bones earlier than X-rays. Whole-body scan takes 3-4 hours after injection.

\section*{Alternative Names}
Nuclear medicine bone scan, bone scintigraphy, technetium-99m bone scan, whole-body bone scan


\section*{Principle}
Bone scan uses technetium-99m-labeled bisphosphonates (MDP or HDP) injected intravenously. The tracer is taken up by bone in proportion to blood flow and osteoblastic activity. A gamma camera detects the emitted 140 keV gamma photons. Increased uptake (hot spots) indicates increased bone turnover from metastases, infection, or fracture.

\section*{History}
Nuclear medicine bone scanning was developed in the 1960s using technetium-99m labeled diphosphonates. It became a routine clinical tool in the 1970s for detecting metastatic bone disease.

\section*{Why This Test Is Ordered}
Your doctor orders a bone scan when they suspect bone metastases (cancer that has spread to your bones), when you have a suspected bone infection, when you have a fracture that does not show up on X-ray, or when evaluating metabolic bone disease. It answers the question: are there areas of abnormal bone activity throughout your skeleton?

\section*{Is This a Primary or Secondary Test}
A secondary test. Usually ordered after other tests suggest a bone problem, or as part of cancer staging.

\section*{How This Test Is Performed}
A small amount of radioactive material is injected into your vein. You wait three to four hours while the material distributes to your bones. Then you lie on a table while a special camera takes images of your entire skeleton. The test takes thirty to sixty minutes. There is no pain involved.

\section*{What Preparation Is Needed}
No special preparation is needed. Tell your doctor if you are pregnant or breastfeeding.

\section*{Contraindications and Precautions}
Bone scans use a small amount of radioactive tracer and are generally not performed during pregnancy. Inform the team if you are pregnant or breastfeeding. The tracer passes into breast milk; discuss timing with your doctor if you are breastfeeding.

\section*{Risks}
A bone scan involves a small dose of radioactive material injected into a vein. The radiation exposure is comparable to a CT scan. The tracer may cause mild allergic reaction at the injection site. Inform the team if you are pregnant or breastfeeding.

\section*{What Happens After the Test}
You can resume normal activities immediately. The radioactive material will naturally decay and leave your body within a day or two. Drink plenty of fluids to help flush it out. Results typically take one to three days.

\section*{Understanding Your Results}

\section*{Normal Results}
A normal result shows symmetric tracer uptake throughout the skeleton with no areas of abnormally increased or decreased activity.
\clearpage
